\section{Bài toán quy hoạch toàn phương}
Quy hoạch toàn phương (Quadratic Programming - QP) là một lớp bài toán tối ưu lồi có vai trò quan trọng trong điều khiển robot, ước lượng tham số, phân phối lực và nhiều ứng dụng liên quan đến bất đẳng thức tuyến tính. QP đặc biệt phù hợp cho các bài toán điều khiển thời gian thực bởi tính chất hội tụ nhanh, nghiệm duy nhất và khả năng mô tả tự nhiên nhiều loại ràng buộc vật lý.Bài toán QP tổng quát được viết dưới dạng:
\begin{equation}
\label{eq:qp_general}
\begin{aligned}
\min_{x \in \mathbb{R}^n} \quad 
    & \frac{1}{2} x^\top H x + c^\top x, \\[4pt]
\text{s.t.} \quad 
    & A x \le b, \\
    & A_{\mathrm{eq}} x = b_{\mathrm{eq}},
\end{aligned}
\end{equation}
trong đó $H \in \mathbb{R}^{n \times n}$, $c \in \mathbb{R}^n$, và các ma trận $A, b, A_{\mathrm{eq}}, b_{\mathrm{eq}}$ mô tả ràng buộc bất đẳng thức và đẳng thức.

Nếu ma trận $H$ là xác định dương ($H > 0$), bề mặt của hàm chi phí trở thành một parabol đa chiều có duy nhất một điểm cực tiểu toàn cục. Điều này không những đảm bảo nghiệm tìm được là nghiệm tối ưu, mà còn loại bỏ hoàn toàn nguy cơ rơi vào các điểm cực trị cục bộ như trong các bài toán phi tuyến. Chính sự đơn giản và chắc chắn này làm cho QP trở thành một trong những công cụ đáng tin cậy nhất trong điều khiển thời gian thực, nơi tính ổn định của nghiệm quan trọng không kém gì tốc độ tính toán. Đây là một trong những đặc điểm quan trọng nhất, mang lại sự ổn định và khả năng dự đoán được của nghiệm tối ưu, đặc biệt trong các ứng dụng điều khiển robot, nơi tín hiệu điều khiển phải thay đổi nhịp nhàng theo thời gian và không được phép dao động thất thường.

Việc kết hợp các ràng buộc tuyến tính dạng $A x \le b$ cho phép bài toán mô tả tự nhiên nhiều giới hạn vật lý của hệ thống như giới hạn lực, giới hạn mô-men. Vì các ràng buộc tuyến tính không phá vỡ tính lồi của bài toán, toàn bộ cấu trúc tối ưu hóa vẫn giữ được tính chất hội tụ tốt và dễ dàng giải bằng các thuật toán hiệu quả. Trong trường hợp tập nghiệm thỏa mãn ràng buộc là không rỗng và hàm mục tiêu lồi chặt, nghiệm tối ưu vừa tồn tại vừa duy nhất, vừa không bị phụ thuộc vào khởi tạo ban đầu. hờ đặc tính đơn giản nhưng mạnh mẽ này, QP trở thành công cụ mô tả và giải quyết nhiều bài toán trong điều khiển robot, đặc biệt những bài toán có hình thức hàm chi phí toàn phương và ràng buộc tuyến tính, như phân phối lực, điều khiển lực-vị trí lai, tối ưu chuyển động và kiểm soát tiếp xúc. Khả năng diễn đạt đầy đủ cấu trúc toán học của hệ thống thông qua QP mà vẫn đảm bảo khả năng tính toán thời gian thực là lý do QP được sử dụng rộng rãi trong nhiều hệ thống robot hiện đại.